\section{Conclusion}
\label{sec:conclusion}

This paper has developed a simple observation into a modest framework. The observation is that the Magic Gem characterization of magic squares---zero covariance between position and value---extends naturally to arbitrary point configurations with real-valued assignments. The framework provides a language for discussing spin-position balance in geometric settings and places magic squares within a broader context.

Our main theoretical result, Theorem~\ref{thm:connection}, establishes that classical magic squares are precisely the balanced (zero-covariance) configurations on a planar lattice with integer values. This is not a deep theorem---the proof requires only careful bookkeeping---but it clarifies the relationship between the combinatorial definition of magic squares and their statistical characterization.

Our numerical experiments explored the spin balance energy landscape for several symmetric geometries. Magic configurations exist for all tested systems with four or more points and can be found by solving a linear system. These configurations are isolated minima (in low-dimensional systems) or form linear subspaces (in higher-dimensional systems). Random configurations typically have substantial energy, and magic configurations are correspondingly rare.

We have been deliberately conservative in our claims. The framework we present is elementary, the results are straightforward, and the applications remain speculative. We have not demonstrated that the covariance perspective provides new insights into any specific problem in physics or combinatorics.

What we have shown is that the ``magic'' property of magic squares is not an isolated curiosity but an instance of a natural geometric condition. Whether this observation leads to useful developments---in the theory of magic squares, in molecular magnetism, or elsewhere---remains to be seen.

We close with a remark on the nature of this work. The connection between magic squares and covariance is, in retrospect, almost obvious: if row sums and column sums are all equal, then position cannot predict value, and covariance measures exactly this prediction. The contribution of this paper is not discovering something hidden but rather making explicit something that was implicit. Sometimes the most useful thing a piece of mathematics can do is provide a clear lens through which to view familiar objects.

\section*{Acknowledgments}

The author thanks the developers of the Magic Gem framework~\cite{magic_gems_2025} for inspiring this work. Computational resources were provided by [Institution]. The author is grateful to [Colleagues] for helpful discussions.

\section*{Data Availability}

Python code implementing all calculations and generating all figures is available in the supplementary materials.

